%% ch10.tex — Chapter 10: Domain Applications: Social Networks,
%%                         Time Series, and Healthcare
%% Part of "Data Engineering" textbook

\chapter{Domain Applications: Social Networks, Time Series, and Healthcare}
\label{ch:domain-apps}

\epigraph{\textit{``Big data hubris is the often implicit assumption that
big data are a substitute for, rather than a supplement to, traditional
data collection and analysis.''}}{David Lazer, Ryan Kennedy, Gary King
\& Alessandro Vespignani, \textit{Science}, 343(6176), 2014}

\vspace{4pt}

\begin{chapterlearning}
After completing this chapter, you will be able to:
\begin{enumerate}
  \item Explain four structural properties of social media graphs---power-law
        degree distribution, small-world diameter, asymmetric directed
        edges, and community structure---and state the implications of each
        for analysis at big data scale.
  \item Describe the Louvain community detection algorithm: its modularity
        objective, two-phase greedy procedure, and $O(n \log n)$ complexity,
        and explain why it is preferred over spectral clustering for
        billion-node graphs.
  \item Define influence maximisation, state Kempe, Kleinberg \&
        Tardos's NP-hardness result and their greedy $(1-1/e)$
        approximation guarantee, and identify which nodes to seed for
        maximum information spread.
  \item Apply VADER sentiment analysis to social media text, interpret
        the compound score, and describe how VADER handles capitalisation,
        punctuation, degree modifiers, and negation differently from
        bag-of-words approaches.
  \item Explain what makes time series data fundamentally different from
        i.i.d.\ tabular data: autocorrelation, seasonality, trend, and
        non-stationarity, and describe how ARIMA$(p,d,q)$ addresses each.
  \item Describe Facebook Prophet's additive decomposition
        $y(t) = g(t) + s(t) + h(t) + \varepsilon(t)$, explain how
        automatic changepoint detection and Fourier seasonality work,
        and explain why the model is preferred over ARIMA for business
        forecasting at scale.
  \item Describe two anomaly detection methods for time series---STL
        decomposition with residual IQR thresholding and Isolation Forest
        ---and explain the circumstances under which each is appropriate.
  \item Design a vectorised batch forecasting pipeline using PySpark's
        \code{applyInPandas} that fits one Prophet model per partition
        across thousands of time series simultaneously.
  \item Identify and contrast four healthcare data modalities---EHR/FHIR,
        genomics (VCF/Parquet), wearables (Kafka streaming), and medical
        imaging (DICOM)---and select the appropriate storage, processing,
        and privacy technique for each.
  \item Explain three privacy-preserving techniques for healthcare analytics
        ---safe harbour de-identification, differential privacy, and
        federated learning---and identify which regulatory framework
        (HIPAA, GDPR, or CCPA) mandates each.
\end{enumerate}
\end{chapterlearning}

\bigskip

The preceding nine chapters built the complete data engineering stack:
distributed storage (Chapter~\ref{ch:hdfs}), batch processing
(Chapter~\ref{ch:mapreduce}), resource management and file formats
(Chapter~\ref{ch:formats}), in-memory processing with Spark
(Chapter~\ref{ch:spark}), SQL at scale (Chapter~\ref{ch:sql-scale}),
streaming pipelines and data lakes (Chapter~\ref{ch:pipelines}),
and visualization (Chapter~\ref{ch:visualization}).
This chapter applies the complete stack to three high-impact application
domains that collectively define where big data systems deliver the most
measurable value: social media analytics, time series forecasting, and
healthcare.

Each domain presents a distinct combination of data characteristics
(graph-structured, temporally ordered, and privacy-sensitive respectively)
and a distinct set of algorithmic techniques. All three domains share the
same underlying infrastructure---HDFS or object storage for persistence,
Kafka for event ingestion, Spark for distributed processing, and Grafana
or Superset for visualization---demonstrating that the engineering stack
built in previous chapters is not domain-specific but universally applicable.

The chapter closes with three research spotlights and a case study
that illustrates both the power and the peril of big data in public
health: the rise and fall of Google Flu Trends, and the more robust
Prophet-based syndromic surveillance architecture that replaced it.

%%====================================================================
\section{Three Domains of Big Data Application}
\label{sec:domain-overview}
%%====================================================================

The same infrastructure serves radically different application domains.
Table~\ref{tab:three-domains} summarises the characteristics, tools,
and primary analytical tasks for the three domains covered in this chapter.

\begin{table}[htbp]
\centering
\renewcommand{\arraystretch}{1.4}
\begin{tabular}{lp{3.0cm}p{3.0cm}p{3.2cm}}
  \toprule
  \textbf{Domain} & \textbf{Data Characteristics}
    & \textbf{Primary Algorithms} & \textbf{Domain Tools} \\
  \midrule
  Social Media
    & Graph-structured, temporal, text-heavy, multimedia
    & Louvain, PageRank, influence maximisation, VADER, LSH
    & GraphX, NetworkX, Neo4j, VADER \\
  Time Series
    & Ordered, autocorrelated, seasonal, non-stationary
    & ARIMA, Prophet, LSTM, STL, Isolation Forest
    & statsmodels, Prophet, Spark MLlib \\
  Healthcare
    & Semi-structured EHR, genomic, streaming wearables, binary imaging
    & Random Forest, Gradient Boosting, CNN, AlphaFold
    & FHIR, GATK, TensorFlow, HAPI \\
  \midrule
  \multicolumn{4}{c}{\textit{Shared infrastructure: HDFS / S3 $\cdot$
    Kafka $\cdot$ Spark $\cdot$ Hive $\cdot$ Grafana / Superset}} \\
  \bottomrule
\end{tabular}
\caption{Three big data application domains. All three share the same
         distributed infrastructure; domain-specific tools sit above
         the common stack.}
\label{tab:three-domains}
\end{table}

%%====================================================================
\section{Domain A: Social Media Analytics}
\label{sec:social-media}
%%====================================================================

Social media platforms are among the largest-scale data systems humans
have ever built. Twitter (now X) processes 500 million posts per day;
Facebook logs 100 billion interactions; TikTok serves 1 billion video
views. In many markets, Facebook remains the primary channel for crisis
communication, political discourse, and commercial marketing. The data they generate is
graph-structured, temporally indexed, linguistically complex, and at a
scale that requires the full Hadoop/Spark stack to process.

%--------------------------------------------------------------------
\subsection{Graph Structure and Scale}
\label{subsec:social-graph}
%--------------------------------------------------------------------

Social platforms are fundamentally \term{directed graphs}: users are
nodes; following, friending, retweeting, and replying are edges. The
properties of these graphs differ dramatically from the random graphs
assumed by naive statistical models.

\begin{defbox}{Social Network Graph Properties}
A \textbf{social network graph} $G = (V, E)$ has four structural
properties that distinguish it from random graphs:
\begin{enumerate}
  \item \textbf{Power-law degree distribution}: the fraction of nodes
        with degree $k$ follows $P(k) \propto k^{-\alpha}$ for some
        $\alpha \in (2, 3)$. A small number of ``hub'' nodes (celebrities,
        media accounts) have millions of followers; the vast majority
        have tens. This is called a \emph{scale-free network} after
        Barabási \& Albert (1999).
  \item \textbf{Small-world property}: despite billions of nodes, the
        average shortest path between any two nodes is very short---Kwak
        et al.\ (2010) found 4.12 hops on Twitter's 41.7-million-node
        graph. Milgram's ``six degrees of separation'' (1967) observed
        this empirically in 1967; it is now a mathematical consequence
        of power-law degree distributions.
  \item \textbf{Asymmetric directed edges}: on Twitter, only 22.1\%
        of following relationships are reciprocal. On Facebook (undirected,
        mutual friendship), reciprocity is 100\%. The directedness of the
        graph determines how information propagates.
  \item \textbf{Community structure}: nodes cluster into densely
        interconnected groups (communities) that are sparsely connected
        to other communities. Political echo chambers, professional
        networks, and fan communities are all examples.
\end{enumerate}
\end{defbox}

\textbf{Storage and processing.} A social graph with $N$ nodes and $M$
edges requires $O(N + M)$ storage. Twitter's 2010 graph had
$N = 41.7$~million users and $M = 1.47$~billion directed edges,
storing roughly 40~GB of compressed adjacency data. At 2024 scale
($N = 500$~million users), the full following graph is $\sim$500~GB
compressed. Storing and querying this graph requires distributed systems:
\begin{itemize}
  \item \textbf{Apache GraphX} (part of Spark) for large-scale iterative
        graph algorithms (PageRank, connected components, triangle count)
        running on the full graph in Spark's in-memory model.
  \item \textbf{Neo4j} for online, low-latency graph queries
        (``show me the mutual friends of user X and user Y'')
        where Spark's batch latency would be unacceptable.
  \item \textbf{Apache Hive / Parquet} for offline analytical queries
        over the graph's edge list (``how many edges were added in Q1 2024?'').
\end{itemize}

%--------------------------------------------------------------------
\subsection{Community Detection: The Louvain Algorithm}
\label{subsec:community}
%--------------------------------------------------------------------

A \term{community} in a graph is a group of nodes that are more densely
connected to each other than to the rest of the graph. Detecting communities
reveals interest groups, political factions, professional clusters, or
coordinated bot networks. It is one of the most practically important
operations in social network analysis.

\begin{defbox}{Modularity and the Louvain Algorithm (Blondel et al., 2008)}
\textbf{Modularity} $Q$ measures the fraction of edges that fall within
communities minus the expected fraction under a null model of random
edge placement:
\[
  Q = \frac{1}{2m} \sum_{i,j}
  \left[ A_{ij} - \frac{k_i k_j}{2m} \right] \delta(c_i, c_j)
\]
where $A_{ij}$ is the adjacency matrix, $k_i$ is node $i$'s degree,
$m$ is the total number of edges, and $\delta(c_i, c_j) = 1$ if nodes
$i$ and $j$ are in the same community.

The \textbf{Louvain algorithm} maximises $Q$ in two phases repeated
iteratively:
\begin{enumerate}
  \item \textbf{Local optimisation:} each node is moved to the neighbour
        community that maximises the modularity gain $\Delta Q$. A node
        stays in its own community if no move improves $Q$. This passes
        over all nodes until no further improvement is possible.
  \item \textbf{Network aggregation:} communities become super-nodes.
        Edges between communities become weighted edges between
        super-nodes. The algorithm is then applied to the aggregated
        network.
\end{enumerate}
Complexity: $O(n \log n)$ on sparse graphs---scalable to graphs with
hundreds of millions of nodes.
\end{defbox}

\begin{lstlisting}[style=pyspark, caption={Community detection with Louvain
  in Python (networkx + python-louvain)}]
import networkx as nx
import community as community_louvain   # pip install python-louvain
import matplotlib
matplotlib.use("Agg")
import matplotlib.pyplot as plt

# Construct a simple directed follow graph (Twitter-style)
# In production: load edge list from HDFS Parquet via Spark
G_directed = nx.DiGraph()
edges = [
    (1, 2), (2, 1), (1, 3), (3, 1),
    (2, 3), (3, 4), (4, 5), (5, 4),
    (4, 6), (6, 5), (5, 7), (7, 5),
    (1, 8), (8, 4),   # bridge nodes
]
G_directed.add_edges_from(edges)

# Louvain requires undirected input
G = G_directed.to_undirected()

# --- Apply Louvain ---
partition = community_louvain.best_partition(G)
modularity = community_louvain.modularity(partition, G)
print(f"Modularity Q = {modularity:.4f}")

communities = {}
for node, comm_id in partition.items():
    communities.setdefault(comm_id, []).append(node)

for comm_id, members in sorted(communities.items()):
    print(f"  Community {comm_id}: nodes {sorted(members)}")

# Visualise
pos    = nx.spring_layout(G, seed=42)
colors = [partition[n] for n in G.nodes()]
fig, ax = plt.subplots(figsize=(8, 6))
nx.draw_networkx(G, pos=pos, node_color=colors,
                 cmap=plt.cm.Set1, node_size=500,
                 font_color="white", ax=ax)
ax.set_title(f"Louvain Communities  (Q = {modularity:.3f})")
ax.axis("off")
plt.tight_layout()
plt.savefig("louvain_communities.png", dpi=150)
print("Saved: louvain_communities.png")
\end{lstlisting}

\begin{systemdesign}{Scaling Community Detection to Billion-Node Graphs}
For graphs with hundreds of millions of nodes, Python's
\texttt{python-louvain} library is too slow. Production implementations use:
\begin{itemize}
  \item \textbf{Apache Spark GraphX}: the Pregel API implements
        iterative vertex-program algorithms including Louvain phases;
        the full graph is partitioned across executors and edges are
        shuffled only when communities cross partition boundaries.
  \item \textbf{Neo4j Graph Data Science (GDS) library}: implements
        Louvain natively with a streaming result API; suitable for
        online, low-latency community refresh.
  \item \textbf{GraphSAGE / Graph Neural Networks} (Hamilton et al.,
        2017): for graphs where node feature attributes (text, image
        embeddings) should inform community assignment, GNNs learn
        community representations end-to-end.
\end{itemize}
For the \emph{full} Twitter graph at 500 million nodes, Louvain in
Spark GraphX requires approximately 3--5 iterations to converge,
with each iteration involving one full shuffle of the edge list.
On a 100-executor Spark cluster with 16~GB RAM per executor, this
completes in $\sim$20--40 minutes.
\end{systemdesign}

%--------------------------------------------------------------------
\subsection{Influence Maximisation}
\label{subsec:influence}
%--------------------------------------------------------------------

Given a social network and a budget of $k$ ``seed'' users, the
\term{influence maximisation} problem asks: which $k$ users should
receive an initial message so that the maximum number of users are
ultimately reached through organic sharing?

This problem has commercial applications (viral marketing campaigns,
product seeding), public health applications (vaccination campaign
targeting), and disinformation applications (coordinated bot networks
that seed false narratives at high-PageRank nodes).

\begin{defbox}{Influence Maximisation: NP-Hard with
  $(1 - 1/e)$ Approximation (Kempe, Kleinberg \& Tardos, 2003)}
Under the \textbf{independent cascade model}, each edge $(u, v)$ has
a probability $p_{uv}$ that user $u$ activates user $v$ in a single time
step. The influence $\sigma(S)$ of a seed set $S$ is the expected
number of activated nodes after cascades run to completion.

\textbf{Problem:} Find $S^* = \arg\max_{|S| \leq k} \sigma(S)$.

\textbf{Hardness:} The problem is NP-hard; $\sigma(S)$ is
\#P-hard to compute exactly.

\textbf{Greedy approximation:} The function $\sigma(\cdot)$ is
\emph{monotone submodular}: adding a node to a larger set provides no
more benefit than adding it to a smaller set. Nemhauser's theorem
guarantees that the greedy algorithm---at each step, add the node that
maximises the marginal gain $\sigma(S \cup \{v\}) - \sigma(S)$---achieves
\[
  \sigma(S_{\text{greedy}}) \geq \left(1 - \frac{1}{e}\right) \sigma(S^*)
  \approx 0.632 \cdot \sigma(S^*)
\]
This is the best polynomial-time approximation ratio possible unless
P = NP.
\end{defbox}

\begin{keyinsight}{PageRank vs.\ Follower Count for Seeding}
Kwak et al.\ (2010) demonstrated that users ranked by \emph{follower count}
and users ranked by \emph{PageRank} are substantially different populations.
A user with 1 million followers who is followed by many high-PageRank users
(e.g., other celebrities and media accounts) has higher influence than a
celebrity with 5 million followers who is followed by low-influence accounts.
For influence maximisation campaigns, seeding at high-PageRank nodes
produces measurably more cascade depth than seeding at high-follower-count
nodes. This is because PageRank correctly captures the \emph{recursive}
nature of influence: being followed by influential people makes you more
influential.
\end{keyinsight}

%--------------------------------------------------------------------
\subsection{Sentiment Analysis: VADER}
\label{subsec:sentiment}
%--------------------------------------------------------------------

Sentiment analysis assigns a polarity (positive, negative, or neutral)
to a piece of text. At social media scale---millions of tweets per
hour---it provides real-time measures of public opinion, brand sentiment,
and crisis emotion.

Standard bag-of-words and TF-IDF approaches fail on social media text
because they ignore the conventions of online communication:
capitalisation for emphasis, punctuation for intensity, emojis for
emotion, and degree modifiers (``very'', ``kind of'', ``absolutely'')
that shift sentiment magnitude.

\begin{defbox}{VADER: Valence Aware Dictionary and sEntiment Reasoner
  (Hutto \& Gilbert, 2014)}
\textbf{VADER} is a lexicon-and-rule-based sentiment analyser specifically
designed for social media text. Its lexicon assigns a sentiment valence
score to 7{,}517 lexical features (words, emojis, slang terms).
Five heuristic rules modify these base scores:
\begin{enumerate}
  \item \textbf{Capitalisation}: \texttt{GREAT} is scored higher than
        \texttt{great} (amplification factor: $+0.733$).
  \item \textbf{Punctuation}: each exclamation mark adds $+0.292$ to
        a positive score up to a maximum of three.
  \item \textbf{Degree modifiers}: ``very good'' $>$ ``good'' $>$
        ``kind of good''. Boosters add $+0.293$; dampeners subtract.
  \item \textbf{Negation}: a negation word within three positions
        of a sentiment word inverts and attenuates the score by
        multiplying by $-0.74$: ``not good'' $\to$ negative.
  \item \textbf{Contrastive conjunctions}: ``but'' signals that the
        clause following it is more salient than the clause preceding.
\end{enumerate}
Output: a \textbf{compound score} $\in [-1, +1]$. Convention:
$> 0.05$ is positive; $< -0.05$ is negative; otherwise neutral.
\end{defbox}

\begin{lstlisting}[style=pyspark, caption={VADER sentiment in a PySpark
  pipeline (UDF on tweet DataFrame)}]
from pyspark.sql import SparkSession
from pyspark.sql.functions import udf, col, window, avg
from pyspark.sql.types import FloatType
from vaderSentiment.vaderSentiment import SentimentIntensityAnalyzer

spark = SparkSession.builder.appName("TwitterSentiment").getOrCreate()
spark.sparkContext.setLogLevel("WARN")

# Broadcast the VADER analyser  --  one object reused per executor
sid_broadcast = spark.sparkContext.broadcast(SentimentIntensityAnalyzer())

@udf(returnType=FloatType())
def vader_compound(text: str) -> float:
    if text is None:
        return 0.0
    sid = sid_broadcast.value
    return float(sid.polarity_scores(text)["compound"])

# Load Twitter stream from Kafka (see ch08 for Kafka source config)
tweets_df = (spark.readStream
    .format("kafka")
    .option("kafka.bootstrap.servers", "kafka:9092")
    .option("subscribe", "twitter-bd-2024")
    .load()
    .selectExpr("CAST(value AS STRING) AS text",
                "timestamp"))

# Apply VADER UDF
sentiment_df = tweets_df.withColumn(
    "compound", vader_compound(col("text")))

# Aggregate: average sentiment per 5-minute tumbling window
agg_df = (sentiment_df
    .groupBy(window(col("timestamp"), "5 minutes"))
    .agg(
        avg("compound").alias("avg_sentiment"),
        count("*").alias("tweet_count")
    ))

# Stream to console (replace with InfluxDB sink for Grafana)
query = (agg_df.writeStream
    .outputMode("complete")
    .format("console")
    .option("truncate", False)
    .start())

query.awaitTermination()
\end{lstlisting}

%--------------------------------------------------------------------
\subsection{Trend Detection and the Twitter News-Media Finding}
\label{subsec:trends}
%--------------------------------------------------------------------

Kwak et al.\ (2010) analysed 106 million tweets and found that 85\%
of Twitter's trending topics were \emph{headline news events}---not
personal conversations or social memes. The median duration of a
trending topic was 17.5 hours, matching the lifespan of a news cycle.
Breaking news topics reached their peak tweet rate within 15 minutes
of the event---faster than most professional news agencies.

This finding established Twitter as a \term{real-time news medium}
rather than a social network, with profound implications for public
health surveillance: if Twitter reflects real-world events within
minutes, then flu-related tweets should correlate with actual flu
incidence faster than physician-reported surveillance data.

\textbf{Trend detection algorithm.} Twitter's trending topics algorithm
detects sudden relative spikes in phrase frequency using a variant of
Locality-Sensitive Hashing (LSH) on character $n$-grams. A phrase that
normally appears 10 times per hour and suddenly appears 10{,}000 times
per hour is ``trending''---regardless of its absolute frequency. This
\emph{relative} thresholding is what distinguishes a trend from simply
a very popular phrase.

\begin{caution}{Ethical Risks of Social Media Analytics}
The same techniques that enable flu surveillance and brand monitoring
enable political surveillance and targeted manipulation:
\begin{itemize}
  \item Sentiment analysis on political posts can be weaponised to
        identify and suppress dissent---flagging negative sentiment
        about government policy for content removal.
  \item Influence maximisation can be used by coordinated inauthentic
        behaviour campaigns (bot networks) to artificially amplify
        misinformation by seeding it at high-PageRank nodes.
  \item Community detection can expose political affiliations, religious
        identities, or sexual orientations of users who have not
        consented to disclose them---violating privacy through
        \emph{inference}.
\end{itemize}
Any deployment of social media analytics must include an explicit
purpose-limitation policy (stating what the analysis will and will not
be used for) and an independent algorithmic audit.
\end{caution}

%%====================================================================
\section{Domain B: Time Series Analysis}
\label{sec:timeseries}
%%====================================================================

Time series data is fundamentally different from i.i.d.\ tabular data.
A transaction table can be analysed by shuffling its rows without loss
of information. A time series cannot: the temporal order of observations
is itself the information. This section covers the properties that make
time series special, the classical ARIMA model, Facebook Prophet's
modern approach to large-scale forecasting, and two anomaly detection methods.

%--------------------------------------------------------------------
\subsection{Properties of Time Series Data}
\label{subsec:ts-properties}
%--------------------------------------------------------------------

\begin{defbox}{Time Series: Five Structural Properties}
A \textbf{time series} $\{y_t\}_{t=1}^{T}$ is a sequence of observations
indexed by time. Five structural properties must be understood before
modelling:
\begin{enumerate}
  \item \textbf{Autocorrelation}: $y_t$ is correlated with $y_{t-1}$,
        $y_{t-2}$, \ldots at multiple lags. Standard regression that
        assumes $\text{Cov}(y_t, y_{t-k}) = 0$ for $k > 0$ is
        invalid---its parameter estimates are unbiased but its standard
        errors and $p$-values are wrong.
  \item \textbf{Trend}: a long-run upward or downward drift.
        Trend must be removed or modelled explicitly; a model that ignores
        a positive trend will persistently under-predict future values.
  \item \textbf{Seasonality}: regular cyclical patterns that repeat
        with a fixed period---daily (web traffic peaks in the evening),
        weekly (retail sales peak on Friday), or annual (flu season peaks
        in January--February). Multiple simultaneous seasonalities are
        common.
  \item \textbf{Non-stationarity}: the mean or variance of the series
        changes over time. Most classical models (ARIMA, regression)
        require \emph{stationarity} (constant mean and variance).
        Non-stationarity is typically removed by \emph{differencing}:
        $\nabla y_t = y_t - y_{t-1}$.
  \item \textbf{Irregularity (Holiday Effects)}: one-off events
        (public holidays, national elections, product launches, natural
        disasters) cause sharp deviations from the seasonal pattern.
        These cannot be captured by a Fourier seasonality model and must
        be encoded explicitly.
\end{enumerate}
\end{defbox}

%--------------------------------------------------------------------
\subsection{Classical Forecasting: ARIMA}
\label{subsec:arima}
%--------------------------------------------------------------------

\begin{defbox}{ARIMA$(p, d, q)$ — AutoRegressive Integrated Moving Average}
The \textbf{ARIMA}$(p, d, q)$ model has three components:
\begin{itemize}
  \item \textbf{AR$(p)$}: the current value is a linear function of the
        past $p$ values: $y_t = \phi_1 y_{t-1} + \cdots + \phi_p y_{t-p} + \varepsilon_t$.
  \item \textbf{I$(d)$}: the series is differenced $d$ times to achieve
        stationarity. First-order differencing ($d=1$) removes linear
        trend; second-order ($d=2$) removes quadratic trend.
  \item \textbf{MA$(q)$}: the current value depends on the past $q$
        forecast errors: $y_t = \varepsilon_t + \theta_1 \varepsilon_{t-1}
        + \cdots + \theta_q \varepsilon_{t-q}$.
\end{itemize}
Parameter selection uses the \textbf{Akaike Information Criterion (AIC)};
\texttt{auto.arima} (R) and \texttt{pmdarima.auto\_arima} (Python)
automate this search.

\textbf{When ARIMA is appropriate}: short stationary series with no
strong seasonality, requiring exact confidence intervals. Small
datasets ($T < 500$) where Prophet's Bayesian MAP estimation overfits.
Regulatory contexts (financial forecasting) that require documented
statistical assumptions.

\textbf{When ARIMA is \emph{not} appropriate}: multiple overlapping
seasonalities (ARIMA's SARIMA extension handles only one); 10{,}000+
parallel series requiring automated fitting without expert tuning of
$p$, $d$, $q$ for each; irregular calendar effects (holidays, special
events) that fall outside the seasonal period.
\end{defbox}

%--------------------------------------------------------------------
\subsection{Facebook Prophet: Decomposable Forecasting at Scale}
\label{subsec:prophet}
%--------------------------------------------------------------------

Prophet \autocite{Prophet2018} was developed at Facebook's Core Data
Science team to address a concrete business need: daily forecasts for
thousands of internal metrics (daily active users, ad impressions,
revenue) that had to be produced automatically, without a statistician
tuning ARIMA parameters for each series.

\begin{defbox}{Prophet's Additive Decomposition (Taylor \& Letham, 2018)}
Prophet models a time series as an additive composition of four
interpretable components:
\[
  y(t) = g(t) + s(t) + h(t) + \varepsilon(t)
\]
\begin{itemize}
  \item \textbf{Trend} $g(t)$: piecewise linear (or logistic) growth.
        Structural trend changes (\emph{changepoints}) are detected
        automatically by placing potential changepoints at regular
        intervals and imposing a sparse prior
        $\delta_j \sim \mathrm{Laplace}(0, \tau)$ on the slope changes.
        The hyperparameter $\tau$ (called \texttt{changepoint\_prior\_scale})
        controls trend flexibility: small values produce smooth trends;
        large values produce flexible trends that may overfit.
  \item \textbf{Seasonality} $s(t)$: a Fourier series captures periodic
        patterns at multiple resolutions simultaneously:
        \[
          s(t) = \sum_{n=1}^{N}
          \left(a_n \cos\frac{2\pi n t}{P} +
                b_n \sin\frac{2\pi n t}{P}\right)
        \]
        where $P$ is the period (365.25 for annual; 7 for weekly).
        Multiple periods (annual \emph{and} weekly) are summed additively.
        The order $N$ controls smoothness.
  \item \textbf{Holiday effects} $h(t)$: a user-specified table of
        one-off events (Black Friday, national elections, sporting finals,
        bank holidays) is encoded as indicator variables that absorb the
        effect of each event. Events with multi-day impact (e.g., a
        long weekend) are represented with a window.
  \item \textbf{Noise} $\varepsilon(t)$: assumed i.i.d.\ Gaussian;
        the residual not captured by the three structural components.
\end{itemize}
Parameters are estimated via maximum a posteriori (MAP) estimation
using the probabilistic programming language \textbf{Stan}, which is
fast enough to fit thousands of series in parallel.
\end{defbox}

\begin{lstlisting}[style=pyspark, caption={Prophet on one time series:
  e-commerce daily sales volume with Black Friday holiday effect}]
import pandas as pd
import numpy as np
from prophet import Prophet
import matplotlib
matplotlib.use("Agg")
import matplotlib.pyplot as plt

rng = np.random.default_rng(2024)

# Simulate 3 years of daily e-commerce sales volume
dates = pd.date_range("2021-01-01", periods=3*365, freq="D")
t     = np.arange(len(dates))

trend      = 1_000_000 + 500 * t                   # growing trend
weekly_s   = 200_000 * np.sin(2*np.pi * t / 7)    # weekly seasonality
annual_s   = 500_000 * np.sin(2*np.pi * t / 365 - np.pi/2)
noise      = rng.normal(0, 80_000, len(t))
y          = trend + weekly_s + annual_s + noise

df = pd.DataFrame({"ds": dates, "y": y})

# Define shopping holidays (volume spikes during Black Friday / Cyber Monday)
holidays = pd.DataFrame({
    "holiday":      "BlackFriday",
    "ds":           pd.to_datetime([
                        "2021-11-26", "2021-11-29",
                        "2022-11-25", "2022-11-28",
                        "2023-11-24", "2023-11-27",
                    ]),
    "lower_window": -1,   # 1 day before
    "upper_window":  2,   # 2 days after
})

m = Prophet(
    yearly_seasonality  = True,
    weekly_seasonality  = True,
    daily_seasonality   = False,
    holidays            = holidays,
    changepoint_prior_scale = 0.05,    # conservative trend
    seasonality_prior_scale = 10,
)
m.fit(df)

# Forecast 90 days ahead
future   = m.make_future_dataframe(periods=90)
forecast = m.predict(future)

fig1 = m.plot(forecast)
fig1.axes[0].set_title("E-commerce Daily Sales Volume  --  Prophet Forecast")
fig1.savefig("sales_forecast.png", dpi=150, bbox_inches="tight")

fig2 = m.plot_components(forecast)
fig2.savefig("sales_components.png", dpi=150, bbox_inches="tight")
print("Saved: sales_forecast.png and sales_components.png")
\end{lstlisting}

%--------------------------------------------------------------------
\subsection{Anomaly Detection in Time Series}
\label{subsec:anomaly}
%--------------------------------------------------------------------

An \term{anomaly} in a time series is an observation that does not conform
to the expected pattern---either a transient spike, a structural break,
or a slow drift outside expected bounds. Anomaly detection is a critical
component of fraud monitoring, network operations, and predictive
maintenance.

\textbf{Method 1: STL Decomposition.} STL (Seasonal-Trend decomposition
using LOESS) decomposes the time series into trend, seasonal, and
residual components using locally weighted regression (LOESS) for each.
Residuals that exceed $k \times \text{IQR}$ of all residuals are flagged
as anomalies. STL is interpretable (the analyst can inspect each component),
handles missing data gracefully, and works well when the seasonal pattern
is stable.

\textbf{Method 2: Isolation Forest.} An \term{Isolation Forest} is an
ensemble of random trees that \emph{isolates} anomalies by recursively
splitting the feature space at randomly chosen thresholds. An anomaly
requires fewer splits to isolate than a normal observation, because it
lies in a sparse region of the feature space. The anomaly score is the
average \emph{path length} across all trees: short path length indicates
an anomaly. Isolation Forest makes no distributional assumptions, handles
high-dimensional feature spaces, and scales to millions of observations
in Spark MLlib.

\begin{lstlisting}[style=pyspark, caption={Anomaly detection with Isolation
  Forest in PySpark MLlib}]
from pyspark.sql import SparkSession
from pyspark.sql.functions import col, hour, dayofweek, lag
from pyspark.sql.window import Window
from pyspark.ml.feature import VectorAssembler
from pyspark.ml.classification import RandomForestClassifier

# NOTE: Spark MLlib does not have a native IsolationForest (as of Spark 3.5).
# Use the spark-iforest package or scikit-learn on Pandas partitions.
# Here we show the applyInPandas approach with sklearn.

from pyspark.sql.types import StructType, StructField, DoubleType, IntegerType
from sklearn.ensemble import IsolationForest
import pandas as pd

spark = SparkSession.builder.appName("AnomalyDetection").getOrCreate()

# Example: payment agent hourly transaction counts
df = spark.read.parquet("hdfs:///data/payments/agent_hourly_txn/")

# Feature engineering: add lag features
w = Window.partitionBy("agent_id").orderBy("hour_ts")
df = (df
    .withColumn("lag1",  lag("txn_count", 1).over(w))
    .withColumn("lag24", lag("txn_count", 24).over(w))
    .withColumn("hour_of_day", hour(col("hour_ts")))
    .withColumn("dow",         dayofweek(col("hour_ts")))
    .na.drop())

# Schema for the output of our Pandas UDF
output_schema = StructType([
    StructField("agent_id",  df.schema["agent_id"].dataType, True),
    StructField("hour_ts",   df.schema["hour_ts"].dataType,  True),
    StructField("txn_count", DoubleType(), True),
    StructField("anomaly_score", DoubleType(), True),
    StructField("is_anomaly", IntegerType(), True),
])

def detect_anomalies(pdf: pd.DataFrame) -> pd.DataFrame:
    """Fit Isolation Forest per agent partition."""
    features = ["txn_count", "lag1", "lag24", "hour_of_day", "dow"]
    X = pdf[features].fillna(0).values

    if len(X) < 10:          # too few observations
        pdf["anomaly_score"] = 0.0
        pdf["is_anomaly"]    = 0
        return pdf[["agent_id", "hour_ts", "txn_count",
                    "anomaly_score", "is_anomaly"]]

    clf = IsolationForest(n_estimators=100,
                          contamination=0.02,   # expect 2% anomalies
                          random_state=42)
    clf.fit(X)

    pdf["anomaly_score"] = -clf.score_samples(X)    # higher = more anomalous
    pdf["is_anomaly"]    = (clf.predict(X) == -1).astype(int)
    return pdf[["agent_id", "hour_ts", "txn_count",
                "anomaly_score", "is_anomaly"]]

# One IsolationForest per agent; runs in parallel across Spark partitions
anomalies = (df.groupBy("agent_id")
               .applyInPandas(detect_anomalies, schema=output_schema))

anomalies.filter(col("is_anomaly") == 1) \
         .orderBy("anomaly_score", ascending=False) \
         .show(20, truncate=False)
\end{lstlisting}

%--------------------------------------------------------------------
\subsection{Vectorised Batch Forecasting in Spark}
\label{subsec:spark-prophet}
%--------------------------------------------------------------------

A retail chain needs 90-day demand forecasts for 5{,}000 store--product
SKU combinations, each with its own seasonal pattern and promotional
calendar. Fitting 5{,}000 separate Prophet models sequentially on a
single machine would take 8+ hours. Fitting them in parallel across a
Spark cluster reduces this to $\sim$12 minutes on a 20-executor cluster
---a $40\times$ speedup.

The pattern is called \term{vectorised batch forecasting}: each SKU's
time series is assigned to one Spark partition; a Pandas UDF
(\code{applyInPandas}) runs the full Prophet fit-and-predict cycle
inside that partition; results are written as Parquet for serving.

\begin{lstlisting}[style=pyspark, caption={Vectorised Prophet forecasting:
  5\,000 SKUs in parallel with applyInPandas}]
from pyspark.sql import SparkSession
from pyspark.sql.functions import col
from pyspark.sql.types import (StructType, StructField,
    StringType, DateType, DoubleType)
from prophet import Prophet
import pandas as pd
import time

spark = SparkSession.builder \
    .appName("VectorisedProphet") \
    .config("spark.sql.execution.arrow.pyspark.enabled", "true") \
    .getOrCreate()

# Load historical sales: schema (sku_id, date, units_sold)
sales = spark.read.parquet("hdfs:///data/retail/daily_sales/")

# Output schema: what each partition will return
forecast_schema = StructType([
    StructField("sku_id",     StringType(), True),
    StructField("ds",         DateType(),   True),
    StructField("yhat",       DoubleType(), True),
    StructField("yhat_lower", DoubleType(), True),
    StructField("yhat_upper", DoubleType(), True),
])

def fit_and_forecast(pdf: pd.DataFrame) -> pd.DataFrame:
    """Fit Prophet on one SKU's history; return 90-day forecast."""
    sku_id = pdf["sku_id"].iloc[0]

    # Prophet requires columns named 'ds' and 'y'
    train = pdf.rename(columns={"date": "ds", "units_sold": "y"})
    train["ds"] = pd.to_datetime(train["ds"])
    train = train[["ds", "y"]].sort_values("ds")

    if len(train) < 14:     # need at least 2 weeks of data
        return pd.DataFrame(columns=["sku_id", "ds", "yhat",
                                     "yhat_lower", "yhat_upper"])

    m = Prophet(
        yearly_seasonality       = True,
        weekly_seasonality       = True,
        changepoint_prior_scale  = 0.05,
        seasonality_prior_scale  = 10.0,
    )
    m.fit(train)

    future   = m.make_future_dataframe(periods=90)
    forecast = m.predict(future)

    result = forecast[["ds", "yhat", "yhat_lower", "yhat_upper"]].copy()
    result["sku_id"] = sku_id
    result["ds"]     = result["ds"].dt.date
    return result[["sku_id", "ds", "yhat", "yhat_lower", "yhat_upper"]]

t0 = time.perf_counter()

forecasts = (sales
    .repartition(200, "sku_id")   # one or more SKUs per partition
    .groupBy("sku_id")
    .applyInPandas(fit_and_forecast, schema=forecast_schema))

# Write forecasts to Parquet for the inventory management dashboard
(forecasts.write
    .mode("overwrite")
    .partitionBy("sku_id")
    .parquet("hdfs:///data/retail/sku_forecasts_90d/"))

forecasts.count()     # trigger execution
elapsed = time.perf_counter() - t0
print(f"5,000 SKU forecasts completed in {elapsed/60:.1f} minutes")
\end{lstlisting}

%%====================================================================
\section{Domain C: Healthcare Big Data}
\label{sec:healthcare}
%%====================================================================

Healthcare is the domain where big data analytics is most consequential
and most constrained. The potential benefits---earlier disease detection,
personalised treatment, pandemic preparedness---are matched by the
sensitivity of the data and the severity of the consequences of errors.

%--------------------------------------------------------------------
\subsection{Healthcare Data Modalities}
\label{subsec:health-modalities}
%--------------------------------------------------------------------

Healthcare generates four structurally distinct data modalities, each
requiring different storage, processing, and privacy approaches.

\textbf{1. Electronic Health Records (EHR).} EHR systems contain structured
data (ICD-10 diagnosis codes, medication orders, laboratory results,
vital signs), semi-structured data (HL7 FHIR resources, as covered in
Chapter~\ref{ch:integration}), and unstructured data (clinical notes
written by physicians in natural language). A typical hospital generates
2--5~GB of new EHR data per day. At the national level, large healthcare
systems such as the NHS (UK) or Kaiser Permanente (US) hold 15+ years of
digitised records, representing sufficient scale for population-level clinical analytics.

\textbf{2. Genomics.} A single whole-genome sequence at 30$\times$ coverage
produces $\sim$200~GB of raw sequencing reads (FASTQ format), which compress
to $\sim$50~GB as an aligned BAM file and $\sim$1~GB as a called variant VCF
file. A clinical genomics study of 10{,}000 patients generates 2~PB of raw
data---requiring the HDFS-class distributed storage introduced in
Chapter~\ref{ch:hdfs} and the GATK (Genome Analysis Toolkit) pipeline
ported to Spark for distributed variant calling.

\textbf{3. Wearables and IoT.} Continuous monitoring devices (smartwatches,
continuous glucose monitors, cardiac monitors, hospital bedside vital-sign
monitors) generate time-series data at per-second resolution. A single
ICU patient generates $\sim$1~GB of monitoring data per day across
all bedside devices. This is precisely the streaming ingestion scenario
from Chapter~\ref{ch:pipelines}: Kafka receives the device events,
Spark Structured Streaming applies anomaly detection in real time, and
an InfluxDB time-series database stores the per-second aggregates for
Grafana dashboards on the nursing station.

\textbf{4. Medical Imaging.} CT scans (200--500~MB per study), MRI
scans (200~MB--2~GB), digital pathology slides (2--10~GB per slide),
and chest X-rays (50~MB) are stored in \term{DICOM} format---a binary
file format with embedded metadata headers. DICOM is fundamentally
unstructured: there is no predefined row-column schema that SQL can
query. Analysis requires deep learning models (convolutional neural
networks for radiology, vision transformers for pathology slides)
running on GPU infrastructure. PACS (Picture Archiving and Communication
Systems) store these files; the Hadoop stack is not the primary
processing environment for imaging analytics.

\begin{table}[htbp]
\centering
\renewcommand{\arraystretch}{1.4}
\begin{tabular}{lp{2.5cm}p{2.5cm}p{2.5cm}p{2.0cm}}
  \toprule
  \textbf{Modality} & \textbf{Format}
    & \textbf{Storage} & \textbf{Processing} & \textbf{Privacy Risk} \\
  \midrule
  EHR / FHIR  & HL7 FHIR JSON, relational
    & Hive / PostgreSQL & Spark SQL, NLP & High \\
  Genomics     & VCF / BAM / FASTQ
    & HDFS Parquet & GATK on Spark & Very high \\
  Wearables    & JSON events (Kafka)
    & InfluxDB, HDFS & Spark Streaming & Medium \\
  Medical Imaging & DICOM binary
    & PACS / S3 & CNN / ViT on GPU & High \\
  \bottomrule
\end{tabular}
\caption{Healthcare data modalities and their infrastructure requirements.
         Genomic data carries the highest privacy risk: a person's genome
         is permanently identifying and cannot be de-identified by removing
         names and dates.}
\label{tab:health-modalities}
\end{table}

%--------------------------------------------------------------------
\subsection{Clinical Prediction Models}
\label{subsec:clinical-pred}
%--------------------------------------------------------------------

\autocite{Obermeyer2016} argued in a landmark New England Journal of
Medicine paper that machine learning models trained on EHR data at
scale could systematically outperform physician heuristics for three
critical clinical tasks.

\textbf{ICU mortality prediction.} A random forest trained on vital
signs, laboratory values (serum lactate, creatinine, bilirubin,
platelet count), and prior diagnosis codes achieves an AUROC greater
than 0.90 for predicting which ICU patients will die within 24 hours.
Conventional ICU scoring systems (APACHE-IV, SOFA) achieve AUROC
$\approx 0.74$--$0.82$. The difference in AUROC translates to 6--8 hours
of earlier identification of deteriorating patients---time in which
palliative care discussions can occur or escalation of treatment can be
initiated.

\textbf{30-day hospital readmission.} Gradient boosting on structured
discharge records (diagnosis codes, procedure codes, length of stay,
laboratory trends, medication count) identifies high-risk patients for
post-discharge follow-up interventions. When these interventions are
implemented, readmission rates fall by 15--20\%, reducing both patient
suffering and hospital readmission penalties under value-based care
contracts.

\textbf{Sepsis early warning.} Sepsis is a life-threatening organ
dysfunction caused by infection; each hour of delayed treatment increases
mortality by 7\%. A continuous streaming model fed by Kafka (receiving
bedside vital sign updates every 60 seconds) and Spark Structured
Streaming can detect sepsis onset 6 hours before the standard clinical
criteria (SIRS or qSOFA) are met, enabling prophylactic antibiotic
administration.

\begin{keyinsight}{Data Scale Requirements for Clinical ML}
Obermeyer and Emanuel's meta-analysis found that clinical prediction
models trained on fewer than 10{,}000 patient records are generally
unreliable for clinical deployment: the models overfit to local
population characteristics and fail to generalise to patients with
comorbidities not well-represented in the training data. Reliable
generalisation requires 100{,}000+ patient records for common diagnoses
(heart failure, pneumonia, sepsis) and 1{,}000{,}000+ for rare conditions.
Large national health systems (NHS England, Kaiser Permanente, VA Health
System) record tens of millions of inpatient encounters per year---if
integrated and de-identified, sufficient for mortality and readmission modelling.
\end{keyinsight}

%--------------------------------------------------------------------
\subsection{Privacy: HIPAA, GDPR, and the CCPA}
\label{subsec:health-privacy}
%--------------------------------------------------------------------

Healthcare data is the most sensitive personal data category recognised
by law. Three regulatory frameworks govern its handling.

\begin{defbox}{Healthcare Privacy Regulatory Frameworks}
\begin{itemize}
  \item \textbf{HIPAA (USA)} — Health Insurance Portability and
        Accountability Act (1996): requires de-identification of 18
        specific ``protected health information'' (PHI) identifiers
        (name, date of birth, National ID number, zip code, telephone,
        email, account numbers, etc.) before research use. Two de-identification
        methods are defined: the \emph{Safe Harbour} method (remove all 18
        identifiers) and the \emph{Expert Determination} method (a qualified
        statistician certifies that re-identification risk is very small).
  \item \textbf{GDPR (EU)} — General Data Protection Regulation (2018):
        treats health data as a ``special category'' requiring explicit
        consent for processing. Data minimisation (collect only what is
        needed), purpose limitation (use only for the stated purpose), and
        right to erasure (the ``right to be forgotten'') apply. GDPR's
        $\in$20 million or 4\% of global revenue fines have made it the
        effective global standard.
  \item \textbf{CCPA (California)} — California Consumer Privacy Act
        (2020): grants residents the right to know what personal data is
        collected, to opt out of its sale, and to request deletion.
        Health data collected outside HIPAA-covered entities (e.g., by
        consumer wellness apps) falls under CCPA rather than HIPAA.
        Many US states are enacting similar legislation, and CCPA
        serves as a model for comprehensive state-level privacy law.
\end{itemize}
\end{defbox}

\textbf{Technical mitigation techniques.} Four engineering approaches
exist for enabling analytics on sensitive healthcare data without
exposing individual records.

\textbf{1. Safe harbour de-identification.} Remove all 18 HIPAA-specified
identifiers. After safe harbour, the data is no longer legally PHI and
can be shared for research without patient consent. Limitation: dates
must be generalised to year only; zip codes must be generalised to the
first 3 digits---changes that reduce the data's utility for geospatial
and temporal analyses.

\textbf{2. Differential privacy.} Add calibrated Gaussian or Laplace
noise to the output of any aggregate query such that no individual record
can be inferred from the query result. The privacy guarantee is
formal: a mechanism $\mathcal{M}$ is $(\varepsilon, \delta)$-differentially
private if adding or removing any single individual's record changes the
output distribution by at most a factor $e^\varepsilon$ with probability
$1 - \delta$. Smaller $\varepsilon$ (stronger privacy) requires more noise
(less accuracy). Apple, Google, and the US Census Bureau use differential
privacy for aggregate statistics release.

\textbf{3. Federated learning.} Rather than centralising patient data
in a single data lake, each hospital trains the model locally on its
own data. Only the model's gradient updates (parameter changes, not raw
data) are sent to a central coordinator, which aggregates them using
\emph{federated averaging} (McMahan et al., 2017). The raw patient
records never leave the hospital. Federated learning is the standard
approach for multi-site clinical trials and cross-hospital model training
in GDPR-regulated environments.

\textbf{4. Synthetic data generation.} Generative models (Variational
Autoencoders, Generative Adversarial Networks) trained on real patient
records produce \emph{synthetic} patient records that preserve the
statistical properties of the original data while containing no real
individuals. Synthetic data can be shared publicly without legal
restriction. Limitation: the generative model itself must be trained
on real data, so data access remains necessary at the training stage.

%--------------------------------------------------------------------
\subsection{AlphaFold: Big Data and Deep Learning in Science}
\label{subsec:alphafold}
%--------------------------------------------------------------------

AlphaFold~\autocite{AlphaFold2021} is the most dramatic example of
big data combined with deep learning transforming an experimental
science. For 50 years, predicting a protein's three-dimensional structure
from its amino acid sequence was considered one of biology's hardest
problems---the ``protein folding problem.'' Experimental determination
by X-ray crystallography or cryo-electron microscopy takes months per
protein and costs thousands of dollars. As of 2020, the Protein Data
Bank (PDB) contained approximately 170{,}000 experimentally determined
structures---a tiny fraction of the $\sim$250 million known protein
sequences in the UniProt database.

\textbf{AlphaFold 2 (DeepMind, 2020)} solved this problem with:
\begin{itemize}
  \item A transformer-based architecture (\emph{Evoformer}) that operates
        jointly on the amino acid sequence and a \emph{multiple sequence
        alignment (MSA)}---a matrix of thousands of evolutionarily related
        sequences from 250 million protein sequences in UniProt.
  \item Training on all 170{,}000 PDB structures using 128 TPU v3 cores
        over several weeks.
  \item Prediction accuracy: median GDT score $>$92 on the CASP14
        benchmark---the first computational method to achieve expert-level
        accuracy.
\end{itemize}

In 2022, DeepMind released the \emph{AlphaFold Protein Structure Database}:
predicted structures for 200 million proteins across all known organisms,
covering effectively the entire known proteome---a resource that took
experimental biology 70 years to partially populate. The database is
freely accessible and has been used in vaccine design, drug discovery, and
fundamental biology research.

\textbf{Why this is a big data problem.} The key input to AlphaFold
is not just the query sequence---it is the MSA. Generating an MSA for
a single query requires searching 250 million sequences (200+ GB compressed),
aligning the top hits, and encoding the result as a 2-D attention matrix.
This search alone takes 3--5 minutes on a high-memory CPU server. Training
required petabyte-scale data access across the PDB and UniRef90. AlphaFold
is in this sense less a breakthrough in neural network architecture than
a breakthrough in what becomes possible when a deep learning model is
trained on a complete and expertly curated scientific dataset at scale.

\begin{techdebt}{AlphaFold's Limitation: Dynamics and Complexes}
AlphaFold predicts \emph{static} protein structures in isolation.
Two fundamental limitations remain for practical drug discovery:
\begin{enumerate}
  \item Many drug targets are \emph{intrinsically disordered proteins}
        (IDPs) that have no fixed structure---they fold only upon binding
        their interaction partner. AlphaFold's confidence scores for IDPs
        are low and should not be trusted as structural predictions.
  \item AlphaFold predicts monomeric (single-chain) structures.
        Many proteins function as \emph{complexes} (multi-chain assemblies).
        AlphaFold-Multimer (2021) and AlphaFold 3 (2024) address this,
        but with lower accuracy than the monomer model.
\end{enumerate}
In production use for drug discovery pipelines, AlphaFold predictions
are treated as \emph{starting hypotheses} for experimental validation,
not as definitive structures.
\end{techdebt}

%%====================================================================
\section{Research Spotlight}
\label{sec:ch10-research}
%%====================================================================

\begin{researchspotlight}{Kwak et al.\ (2010) · Taylor \& Letham (2018)
  · Obermeyer \& Emanuel (2016) --- Social Networks, Time Series, and
  Clinical ML at Scale}

\textbf{Work 1: Kwak, H., Lee, C., Park, H., \& Moon, S. (2010).
``What is Twitter, a Social Network or a News Media?''
\textit{Proceedings of WWW 2010}, pp.~591--600. ACM.}
($>$7{,}000 citations; KAIST.)

The paper's central question was deceptively simple: is Twitter's graph
structure more like a social network (mutual relationships, dense local
communities, information spreading through friendship chains) or a news
medium (one-to-many broadcasting, shallow cascades, content driven by
events rather than social ties)?

To answer it, Kwak, Lee, Park, and Moon crawled Twitter's full public
graph in June 2009---41.7 million users, 1.47 billion directed following
edges---using breadth-first search via the Twitter API over three weeks.
The resulting 40~GB compressed adjacency dataset was one of the first
big data social graph studies.

Three findings reshaped social media analytics:
\begin{enumerate}
  \item \textbf{Twitter is a news medium.} 85\% of trending topics were
        headline news events, not personal conversations. The median trending
        topic lasted 17.5 hours. Breaking news reached peak tweet rate
        within 15 minutes---faster than news wire services.
  \item \textbf{Asymmetric structure.} Only 22.1\% of following
        relationships are reciprocal---far lower than Facebook's 100\%.
        The in-degree distribution follows a power law with exponent
        $\alpha \approx 2.276$: a small number of celebrity accounts
        have millions of followers, while the vast majority have tens.
        The average shortest path across the full 41.7-million-node
        graph is 4.12 hops.
  \item \textbf{PageRank diverges from follower count.} Users ranked by
        PageRank (recursive influence in the directed graph) differ
        substantially from those ranked by raw follower count.
        A user followed by many high-PageRank accounts outranks a
        celebrity followed by low-influence accounts.
\end{enumerate}

\textbf{Impact.} The news-medium finding directly enabled social media
epidemiology: CDC, WHO, and academic groups built influenza, COVID-19,
and mental health surveillance systems on the Twitter graph model.
The directed-cascade model of retweets is the reference network structure
for misinformation propagation studies (Vosoughi, Roy \& Aral, 2018,
\textit{Science}: ``False news spreads faster, farther, and more broadly
than true news''). Twitter's own Trends algorithm redesigns were shaped
by the paper's finding that trending topics are news events.

\medskip
\textbf{Work 2: Taylor, S.~J., \& Letham, B. (2018). ``Forecasting at
Scale.'' \textit{The American Statistician}, 72(1), 37--45.}
($>$1{,}500 citations; Facebook Core Data Science.)

The Prophet paper addresses a practical industrial problem: Facebook
needed daily forecasts for thousands of internal business metrics
(daily active users, ad impressions, revenue per cohort)---automatically,
without a statistician manually tuning ARIMA parameters for each series.
Classical approaches (ARIMA, exponential smoothing) had three failure modes
at this scale: they required expert parameter selection ($p$, $d$, $q$)
that could not be automated reliably across heterogeneous time series;
they struggled with complex multi-period seasonality (annual \emph{and}
weekly simultaneously); and they had no mechanism for handling irregular
calendar effects (holidays that fall on different dates each year).

Prophet's key innovations were:
\begin{enumerate}
  \item \textbf{Interpretable decomposition}: the additive
        $g(t) + s(t) + h(t) + \varepsilon(t)$ structure lets a domain
        analyst (not a statistician) inspect and correct each component
        independently.
  \item \textbf{Automatic changepoint detection}: a sparse Laplace prior
        on trend slope changes detects structural breaks without
        over-segmenting the trend.
  \item \textbf{Fourier seasonality}: annual, weekly, and daily seasonality
        are represented as overlapping Fourier series, capturing all
        simultaneously without the SARIMA workaround of a single fixed period.
  \item \textbf{Stan-based MAP estimation}: fast enough to fit thousands
        of models per day; calibrated uncertainty intervals outperform
        ARIMA's theoretical confidence intervals on real data.
\end{enumerate}

\textbf{Impact.} Prophet became the de facto industry standard for
business time series forecasting within two years: adopted by Airbnb,
Uber, Walmart, and most major tech companies. The \code{applyInPandas}
+ Prophet vectorisation pattern is now the canonical PySpark design
pattern for parallel time series forecasting. In healthcare, Prophet
is used for ICU capacity forecasting, hospital admission nowcasting,
and CDC FluSight flu season projections.

\medskip
\textbf{Work 3: Obermeyer, Z., \& Emanuel, E.~J. (2016). ``Predicting
the Future---Big Data, Machine Learning, and Clinical Medicine.''
\textit{New England Journal of Medicine}, 375(13), 1216--1219.}
($>$2{,}000 citations; NEJM; foundational clinical ML position paper.)

Obermeyer and Emanuel's paper is not a primary research study---it is
a perspective article written for the clinical audience of the NEJM,
arguing on the basis of existing evidence that machine learning models
trained on EHR big data can and should be integrated into clinical
practice. Its influence is exceptional because it appeared in the world's
most prestigious medical journal and was written in language accessible
to practising physicians rather than computer scientists.

Three clinical ML results they synthesised:
\begin{enumerate}
  \item \textbf{ICU mortality}: random forest on vitals and labs achieves
        AUROC $>0.90$ vs.\ APACHE's $\approx 0.74$---earlier identification
        of deteriorating patients.
  \item \textbf{30-day readmission}: gradient boosting on discharge records
        identifies high-risk patients for post-discharge follow-up,
        reducing readmission by 15--20\%.
  \item \textbf{Sepsis}: continuous streaming models detect sepsis 6 hours
        before clinical criteria---time enough for prophylactic antibiotics.
\end{enumerate}

The paper also warned of three failure modes: biased training data (if
certain patient populations are underrepresented, predictions are
unreliable for those groups); overfitting to local population
characteristics that do not generalise to other hospitals; and the
``black-box problem''---clinicians distrust predictions they cannot
explain. These warnings have proved prescient: subsequent audits of
deployed clinical ML models have found racial and socioeconomic bias,
poor cross-site generalisation, and physician distrust leading to
models being ignored in practice even when technically accurate.
\end{researchspotlight}

%%====================================================================
\section{Case Study: CDC Syndromic Surveillance and the Google Flu
  Trends Lesson}
\label{sec:cdc-case}
%%====================================================================

\begin{casestudy}{CDC BioSense and the Prophet-Based Flu Surveillance
  Architecture (2009--2022)}

\textbf{Context: Traditional Surveillance Lag.}
The CDC's traditional influenza surveillance system, ILINet, relies on
$>$3{,}000 sentinel physicians across all 50 US states reporting weekly
counts of influenza-like illness (ILI) patients---defined as fever $> 100$°F
plus cough or sore throat. The data flow is: clinic $\to$ state health
department $\to$ CDC weekly report. The lag between a patient's first
symptoms and the CDC receiving the report is \textbf{1--2 weeks}.
By the time CDC identifies a surge, it may already have passed its peak.

\textbf{Google Flu Trends: The Promise and the Failure.}
In 2008, Ginsberg et al.\ at Google published in \textit{Nature} a
system called Google Flu Trends (GFT): by correlating 45 flu-related
search queries (``flu symptoms'', ``how to treat fever'', etc.) with
ILINet surveillance data, GFT could \emph{predict} flu incidence two weeks
before CDC's official reports with $r^2 = 0.97$ on 5 years of training data.
GFT was cited as the canonical example of ``big data surpassing traditional
surveillance.''

In the 2012--2013 flu season, GFT over-predicted flu incidence by a factor
of \textbf{$2\times$}. Lazer et al.\ (\textit{Science}, 2014)
\autocite{LazerFluTrends2014} diagnosed three causes:
\begin{enumerate}
  \item \textbf{Big Data hubris}: $r^2 = 0.97$ on a 5-year training
        window does not guarantee generalisation when the data-generating
        process is non-stationary. Search behaviour changes as the internet
        population and search interface evolve.
  \item \textbf{Algorithm confounding}: Google's autocomplete algorithm
        (opaque to GFT's developers) suggested flu-related search terms
        to users during flu news coverage, amplifying search volume
        \emph{independently} of actual incidence.
  \item \textbf{Media amplification}: heavy news coverage of a flu season
        drove flu-related searches even among people who were not ill,
        creating a spurious correlation between media attention and
        predicted (but not actual) incidence.
\end{enumerate}
GFT was discontinued as a primary source and eventually taken offline in 2015.

\textbf{BioSense: The Better Architecture.}
The CDC's BioSense Platform represents the mature, post-GFT approach
to syndromic surveillance:
\begin{itemize}
  \item \textbf{Primary signal: direct medical observation.} BioSense
        receives near-real-time electronic health data from $>$3{,}500
        emergency departments, urgent care centres, and hospitals
        via HL7 FHIR messages. Syndromes (influenza-like, gastrointestinal,
        neurological) are automatically extracted from chief complaints
        using NLP. Lag: $<$72 hours vs.\ ILINet's 1--2 weeks.
  \item \textbf{Secondary signal: Twitter/social media.} Twitter flu-related
        keywords are monitored as a \emph{leading indicator} supplementing
        the direct medical signal---not replacing it. The Kwak et al.\ finding
        that Twitter reflects real-world events within minutes makes it
        valuable as a 3--5 day leading signal. But unlike GFT, the Twitter
        signal is an input to a model alongside clinical data, not the
        sole prediction basis.
  \item \textbf{Forecasting: Prophet + Ensemble.}
        CDC's FluSight forecasting challenge (2013--present) pools
        submissions from academic teams worldwide. The top performers
        are always ensemble methods. Prophet is a strong component:
        it models the annual flu season's characteristic double-peaked
        shape (a primary peak in January--February, a secondary peak
        in March) as overlapping seasonal components, and its changepoint
        detection identifies season-start changepoints within 1--2 weeks.
\end{itemize}

\textbf{Scale: 2019--2020 Season.}
\begin{center}
\small
\renewcommand{\arraystretch}{1.25}
\begin{tabular}{ll}
  \toprule
  \textbf{Metric} & \textbf{Value} \\
  \midrule
  Emergency department records processed     & 38 million visits \\
  FHIR messages ingested per day             & $\sim$250{,}000 \\
  HHS regions with independent Prophet models & 60 \\
  Spark executors for Prophet batch inference & 20 \\
  Twitter keyword stream volume monitored    & $\sim$500{,}000 flu-related tweets/day \\
  FluSight ensemble models combined          & 38 (2019--2020 season) \\
  \bottomrule
\end{tabular}
\end{center}

\textbf{Engineering Lessons.}
\begin{enumerate}
  \item \textbf{Proxy data amplifies noise, not signal.} GFT's failure was
        that search queries capture \emph{health anxiety and media attention}
        rather than health incidence. Proxy signals require causal validation
        before policy use; correlation is not sufficient.
  \item \textbf{Algorithm transparency is a public health issue.}
        GFT failed partly because Google's internal autocomplete changes
        were opaque to CDC. When a public health surveillance system
        depends on a third party's algorithm, that algorithm's changes
        become a public health risk.
  \item \textbf{Decomposability improves trust.} Prophet's $g + s + h$
        structure lets epidemiologists inspect each component---an
        epidemiologist can identify which changepoint corresponds to the
        season start, which holiday effect captures school opening, and
        whether the trend component reflects a genuine long-term burden
        change. Black-box deep learning forecasters produce comparable
        accuracy but generate less physician and public health official trust.
  \item \textbf{Ensemble beats any single model.} The CDC FluSight top
        performers from 2013 to 2022 are always ensemble methods.
        No single model---Prophet, ARIMA, LSTM, or human expert---wins
        every season. The value of Prophet in the ensemble is its
        interpretability: when the ensemble disagrees with a Prophet
        component, epidemiologists can inspect the component and understand
        why.
  \item \textbf{Twitter adds lead time, not accuracy.} Social media signals
        are available 3--5 days earlier than BioSense ED data but
        with 30--40\% higher error. The trade-off is latency versus
        precision---an acceptable trade-off for an early warning signal
        but not for a primary surveillance measure.
\end{enumerate}
\end{casestudy}

%%====================================================================
\section*{Chapter Summary}
\label{sec:ch10-summary}
%%====================================================================

Three high-impact application domains demonstrate that the data
engineering stack built in previous chapters is universally applicable.

Social media analytics treats platforms as directed graphs and applies
community detection (Louvain, $O(n \log n)$, modularity maximisation),
influence maximisation (NP-hard, $(1-1/e)$-approximable by greedy),
and real-time sentiment analysis (VADER: compound score handles
capitalisation, punctuation, degree modifiers, and negation) to problems
ranging from viral marketing to public health surveillance. Kwak et al.\ (2010)
established that Twitter is a news medium, not a social network---a finding
that unlocked social media epidemiology.

Time series analysis begins with ARIMA's classical approach (differencing
for stationarity, autoregressive and moving average components for
autocorrelation) and extends to Facebook Prophet's decomposable
$g + s + h$ model (piecewise linear trend with automatic changepoint
detection, Fourier seasonality, holiday tables). Vectorised batch
forecasting via \code{applyInPandas} parallelises one Prophet model
per series across thousands of series simultaneously in Spark. STL and
Isolation Forest address the anomaly detection problem.

Healthcare big data combines four data modalities (EHR/FHIR, genomics,
wearables, DICOM imaging), three privacy frameworks (HIPAA, GDPR, CCPA),
and four privacy-engineering techniques (safe harbour, differential
privacy, federated learning, synthetic data). Clinical prediction models
at scale outperform physician heuristics for ICU mortality, readmission,
and sepsis detection. AlphaFold demonstrated that a 50-year scientific
grand challenge could be solved by training a transformer on a complete
curated database at scale.

The CDC/Google Flu Trends case study provides the chapter's most
important cautionary lesson: high in-sample correlation does not guarantee
out-of-sample accuracy when the data-generating process is non-stationary,
and proxy data captures anxiety and media attention as much as the
phenomenon it purports to measure.

%%====================================================================
\section*{Exercises}
\label{sec:ch10-exercises}
%%====================================================================

\exercisesection{Short Questions}

\sq What are the four structural properties of social network graphs
described in this chapter? For each, state one implication for how
information spreads through the network.

\sq Describe the two phases of the Louvain community detection algorithm.
What is the modularity objective $Q$, and why does maximising it detect
communities?

\sq Why is PageRank a better measure of user influence on Twitter than
raw follower count? Describe a specific scenario where the two measures
would rank users in a different order.

\sq VADER outputs a ``compound score'' for social media text. For each
of the following tweets, state whether the compound score will be higher
or lower than the score for the plain form, and explain which VADER rule
causes the difference:
(a) ``GREAT concert'' vs.\ ``great concert'';
(b) ``not good'' vs.\ ``good'';
(c) ``absolutely fantastic!!!'' vs.\ ``fantastic''.

\sq What is autocorrelation in a time series, and why does it invalidate
standard regression's standard errors? Give an example from a business
time series (e.g., daily e-commerce transactions) where strong autocorrelation
would be expected.

\sq Describe each component of Prophet's model
$y(t) = g(t) + s(t) + h(t) + \varepsilon(t)$.
What hyperparameter controls the flexibility of the trend component,
and what happens if it is set too high?

\sq A Prophet model for weekly hospital outpatient visits shows an
unexplained spike every year in late April and late June. Which Prophet
component is responsible for these spikes? What information would you
add to the model to capture them explicitly?

\sq What are STL decomposition and Isolation Forest? For each, describe
what type of anomaly it is best suited to detect in a time series.

\sq What is the difference between a \emph{safe harbour} de-identification
and \emph{differential privacy} for healthcare data? Give one scenario
where safe harbour is sufficient and one where differential privacy
would also be needed.

\sq What was the core argument of Obermeyer \& Emanuel (2016)? List two
specific clinical tasks where they argued ML models outperform physician
heuristics, and state the metric used to measure outperformance.

\sq Explain why Google Flu Trends over-predicted flu incidence in 2012--2013.
Name all three causes identified by Lazer et al.\ (2014) and describe which
is the most fundamental.

\sq What is \emph{federated learning}? Why is it used in multi-hospital
clinical machine learning, and what is the primary limitation of the
federated averaging algorithm?

\exercisesection{Analytical Problems}

\ap \textbf{Influence Maximisation Calculation.}
A Facebook group for a CDC dengue alert campaign has 1 million users.
The following-graph has a power-law degree distribution with exponent
$\alpha = 2.3$. The campaign budget allows seeding $k = 5$ initial posts.

\begin{enumerate}[label=(\alph*)]
  \item Under the independent cascade model with $p_{uv} = 0.05$ on each
        edge, estimate (qualitatively) whether seeding at the 5 users with
        the highest follower count or the 5 users with the highest PageRank
        will achieve greater spread. Justify your answer using the Kwak
        et al.\ finding on PageRank vs.\ follower count.
  \item The Kempe et al.\ greedy algorithm guarantees a $(1-1/e) \approx 0.632$
        approximation ratio. If the optimal seed set $S^*$ achieves influence
        $\sigma(S^*) = 200{,}000$ users, what is the worst-case influence
        of the greedy seed set? Is this a meaningful guarantee for a
        public health campaign? Discuss.
  \item Kempe et al.'s greedy algorithm requires evaluating $\sigma(S)$
        by Monte Carlo simulation (running the cascade 10{,}000 times).
        For a graph of 1 million nodes, $k = 5$ seed nodes, and 100
        candidate nodes to evaluate at each greedy step, estimate the
        total number of cascade simulations required. Is this tractable
        on a single machine?
\end{enumerate}

\ap \textbf{Prophet Model Diagnostics.}
A Prophet model was fitted to 3 years of weekly dengue case data
for a tropical metropolitan region. The fitted components show:
\begin{itemize}
  \item Trend $g(t)$: an increasing linear trend with two changepoints:
        at week 52 (slope increase) and week 130 (slope decrease).
  \item Seasonality $s(t)$: a strong annual peak centred on weeks 32--38
        (August--September) coinciding with the post-monsoon season.
  \item Residuals $\varepsilon(t)$: mostly near zero except for weeks 88,
        89, and 90, where residuals are $+150$, $+220$, and $+180$ cases
        above the seasonal expected value.
\end{itemize}
\begin{enumerate}[label=(\alph*)]
  \item What real-world events might explain the two changepoints at
        weeks 52 and 130? (Hint: consider vector control interventions,
        climate events, and reporting system changes.)
  \item The large positive residuals in weeks 88--90 were not captured
        by the seasonal component. What type of event would cause a
        3-week anomaly in dengue cases, and how would you modify the
        Prophet model to capture it in future forecasts?
  \item The public health director asks whether the model can predict a
        dengue surge 8 weeks in advance. Describe the specific uncertainty that
        makes 8-week forecasts unreliable in a dengue context (not a
        general statistical argument---specific to dengue epidemiology).
\end{enumerate}

\ap \textbf{Healthcare Privacy Analysis.}
A health system researcher wants to analyse 5 years of inpatient records
(500{,}000 patients) from 10 hospitals to build a 30-day readmission
prediction model. The data includes: patient name, national ID number,
date of birth, home region, ICD-10 diagnosis codes, procedure codes,
admission date, length of stay, and discharge outcome.
\begin{enumerate}[label=(\alph*)]
  \item List all fields that would be removed under HIPAA's safe harbour
        method. Which field is most uniquely identifying and why?
  \item The researcher wants to release a public version of the dataset
        for academic use. Under GDPR's data minimisation principle, which
        fields should be retained and which should be removed or generalised?
        Justify each decision.
  \item A second research team at a partner university wants to train on
        their own 50{,}000-patient dataset and combine insights via
        federation. Why would federated learning be preferable to data
        centralisation? Describe the federated averaging protocol and
        identify one risk.
  \item The final deployed model has AUROC 0.82. The Chief Data Officer
        asks: ``Is the model equally accurate for patients from rural
        regions as for patients from urban centres?'' Describe the
        analysis you would run to answer this question and the corrective
        action if the answer is no.
\end{enumerate}

\ap \textbf{Vectorised Forecasting Performance Analysis.}
A telecom company needs to forecast daily call-drop rates for 6{,}000
cell towers, each with 2 years of daily measurements.
\begin{enumerate}[label=(\alph*)]
  \item Design the Spark job using \code{applyInPandas}. Specify:
        the DataFrame schema, the partition key, the Pandas UDF signature,
        and the output schema for 30-day forecasts.
  \item A single Prophet fit on 730 daily observations takes approximately
        2 seconds. Estimate the wall-clock time to fit all 6{,}000 tower
        models on: (i) a single machine with 8 cores; (ii) a Spark cluster
        with 20 executors each having 8 cores.
  \item The Prophet models for 300 towers (5\%) produce MAPE $> 30\%$
        on a holdout set. Describe three diagnostic checks you would
        perform to identify why these towers have poor forecast accuracy.
  \item Propose a strategy for handling towers with fewer than 60 days
        of data (e.g., new towers installed in the past 2 months).
\end{enumerate}

\exercisesection{Design Problems}

\dpr \textbf{National Disease Surveillance System.}
Design a big data disease surveillance system analogous to the CDC
BioSense platform. The system should detect outbreaks of dengue fever,
influenza, and COVID-19 up to 7 days before the health authority
receives clinical confirmation.

\begin{enumerate}[label=(\alph*)]
  \item \textbf{Data sources}: identify four non-traditional data sources
        analogous to the CDC's syndromic surveillance inputs (social
        media, pharmacy, ER, search). For each, identify the entity that
        controls the data and the access mechanism.
  \item \textbf{Ingestion pipeline}: specify the Kafka topic structure
        (topic name, partition count, key, value format) for each source.
        Explain why Kafka is preferable to a direct database write for
        multi-source signal aggregation.
  \item \textbf{Signal processing}: describe the Spark Structured Streaming
        window aggregation for the social media signal (window size, slide,
        output metric). Describe the Prophet model configuration for the
        weekly hospitalization series.
  \item \textbf{Alert logic}: specify the threshold that triggers an
        outbreak alert. Should the threshold be a fixed case count, a
        percentage above the Prophet forecast, or a combination? Justify.
  \item \textbf{Privacy safeguards}: the social media signal includes
        geotagged posts and the pharmacy signal includes named individuals'
        medication purchases. Specify the de-identification steps before
        these signals enter the pipeline, citing GDPR Article~9 and
        HIPAA Safe Harbour requirements.
\end{enumerate}

\dpr \textbf{Multi-Hospital Clinical ML Pipeline.}
A regional health authority wants to deploy a 30-day readmission
prediction model trained on data from five hospitals in a metropolitan
region. Data sharing between hospitals is restricted by HIPAA / GDPR.

\begin{enumerate}[label=(\alph*)]
  \item Design a federated learning architecture: specify which component
        runs at each hospital, what is transmitted to the central server,
        and how the federated averaging step works.
  \item Specify the minimum feature set that can be computed from ICD-10
        codes, length of stay, and admission month without transmitting
        any PHI identifiers. Explain the trade-off between feature richness
        and privacy risk.
  \item The central model achieves AUROC 0.79 after 10 federated rounds.
        Hospital A's local model (50{,}000 patients) achieves AUROC 0.84.
        The hospital director asks whether to use the global model or the
        local model. Analyse the trade-offs and provide a recommendation.
  \item Describe two bias risks in the federated model arising from
        the dataset heterogeneity across five hospitals (e.g., different
        patient populations, different documentation practices) and
        propose a mitigation for each.
\end{enumerate}

\exercisesection{Programming Exercises}

\pe \textbf{Social Network Analysis with NetworkX and Louvain.}
Construct a social network from a synthetic dataset of 1{,}000 Twitter
users and 15{,}000 directed following edges, then analyse its structure.
\begin{enumerate}[label=(\alph*)]
  \item Generate the synthetic directed graph using \texttt{networkx}
        with a preferential attachment model (\texttt{barabasi\_albert\_graph}).
        Compute: number of nodes, number of edges, average in-degree,
        maximum in-degree, reciprocity rate, and average shortest path
        length (on the largest connected component).
  \item Apply the Louvain algorithm (\texttt{python-louvain}). Report:
        number of communities detected, modularity $Q$, and the size
        of the five largest communities.
  \item Compute PageRank for all nodes. Identify the top-5 nodes by
        PageRank and the top-5 by in-degree. Are they the same nodes?
        Calculate the Spearman rank correlation between PageRank and
        in-degree rankings.
  \item Create a visualisation: draw the network with nodes coloured
        by community assignment (Louvain partition), node size
        proportional to PageRank, and inter-community edges drawn
        as dashed grey lines.
\end{enumerate}

\pe \textbf{Prophet Forecasting with Holiday Effects.}
Using the Prophet library, build a demand forecasting model for an
international retailer with strong year-end seasonal spikes.
\begin{enumerate}[label=(\alph*)]
  \item Generate 3 years of synthetic weekly sales data with: an upward
        linear trend, strong annual seasonality (peak in late November /
        December holiday season), moderate weekly seasonality (peak on
        Saturday), and random Gaussian noise.
  \item Fit two Prophet models: one without holidays and one with
        Black Friday, Cyber Monday, and Christmas dates explicitly encoded
        with a window of $[-3, +3]$ days. Compare the MAPE on a 12-week
        holdout set.
  \item Plot the component decomposition for the model with holidays.
        Identify the estimated holiday effect size (in sales units) for
        Black Friday.
  \item Change \texttt{changepoint\_prior\_scale} from 0.01 to 0.5 in
        steps. Plot how the fitted trend changes. At what value does the
        model begin to overfit to noise in the training data?
\end{enumerate}

\pe \textbf{Healthcare Data Pipeline (FHIR + Spark).}
Build a Spark pipeline that simulates processing EHR data from multiple
hospitals for a 30-day readmission prediction task.
\begin{enumerate}[label=(\alph*)]
  \item Generate 100{,}000 synthetic patient records with fields:
        \texttt{patient\_id} (UUID), \texttt{admission\_date},
        \texttt{region} (50 values), \texttt{primary\_icd10} (50 codes),
        \texttt{los\_days}, \texttt{readmitted\_30d} (binary label,
        15\% positive rate).
  \item Apply safe harbour de-identification: generate a surrogate
        \texttt{anon\_id} (SHA-256 hash of patient UUID), generalise
        \texttt{admission\_date} to \texttt{admission\_year\_month},
        and group \texttt{region} into 10 macro-regions.
  \item Using Spark MLlib, train a Gradient Boosted Tree classifier
        on the de-identified features. Report AUROC, precision, recall,
        and F1 on a 20\% holdout set.
  \item Stratify the AUROC by region. Identify the three regions with
        the lowest AUROC and hypothesise (in 2--3 sentences) why the
        model performs worse in those regions.
\end{enumerate}

%%====================================================================
\section*{Further Reading}
\label{sec:ch10-further}
%%====================================================================

\begin{itemize}
  \item \textbf{\cite{Kwak2010}} --- Five pages; read it in full.
        The network topology findings (Sections 3--4) and the trending
        topics analysis (Section 5) are the most cited. The data collection
        section (Section 2) is worth reading for its description of
        rate-limit-aware BFS crawling---a practical lesson for any
        large-scale API data collection.
  \item \textbf{\cite{Prophet2018}} --- Open access at
        \texttt{peerj.com/preprints/3190}. Read Sections 2 (model)
        and 3 (analyst-in-the-loop). The Appendix contains the Stan
        code for the full model---useful if you want to extend Prophet's
        priors for a specific domain.
  \item \textbf{\cite{Obermeyer2016}} --- Two pages in NEJM.
        Read alongside Caruana et al.'s 2015 KDD paper on ``Intelligible
        Models for Healthcare'' for the interpretability counterargument:
        high-AUROC models that physicians cannot explain are frequently
        ignored in clinical practice.
  \item \textbf{\cite{LazerFluTrends2014}} --- Three pages in
        \textit{Science}. A masterclass in the epistemology of big data:
        why $r^2 = 0.97$ on training data means almost nothing, and why
        algorithm opacity is a scientific problem.
  \item \textbf{\cite{AlphaFold2021}} --- Read the introduction and
        the ``Results'' first paragraph. The Evoformer architecture
        section (Methods) is deep and requires familiarity with
        attention mechanisms; save it for after studying transformers.
  \item \textbf{Leskovec, Rajaraman \& Ullman.
        \textit{Mining of Massive Datasets}, 3rd ed.\ Chapter 10}:
        social network structure, community detection, influence
        propagation. Freely available at \texttt{mmds.org}.
  \item \textbf{Hyndman, R.~J., \& Athanasopoulos, G. (2021).
        \textit{Forecasting: Principles and Practice}, 3rd ed.}
        Freely available at \texttt{otexts.com/fpp3}. The definitive
        modern time series forecasting textbook; Chapters 9--10 cover
        ARIMA and Prophet in depth.
\end{itemize}
